\documentclass[10pt,a4paper]{amsart}
\usepackage[margin=19mm]{geometry}
\usepackage{amsmath,amssymb,mathtools}
\usepackage[scr=rsfs,cal=euler]{mathalfa}
\usepackage{xfrac}
\usepackage[unicode,hidelinks,bookmarks=false]{hyperref}
\newcommand{\ie}{i.e.\ }
\newcommand{\cf}{cf.\ }
\newcommand{\resp}{respectively\ }
\newcommand{\Subsection}{\subsection}
\providecommand{\nobreakdash}{\nobreak\hbox{-}}
\setlength{\emergencystretch}{2em}
\multlinegap0pt
\usepackage{bm}
\usepackage{bbm}
\usepackage{dsfont}
\usepackage{xspace}
\usepackage{cancel}
\usepackage{mathtools}
\usepackage{amsthm}
\makeatletter
\@ifundefined{thm}{\newtheorem{thm}{Theorem}}{}
\makeatother
\makeatletter
\expandafter\ifx\csname varthm\endcsname\relax
\newenvironment{varthm}[1]{\def\varthmname{#1}\begin{thevarthm}}{\end{thevarthm}\def\varthmname{}}
\fi
\makeatother
\title[On combinatorics of plus-one generated line arrangements]{On combinatorics of plus-one generated line arrangements}
\author{Artur Bromboszcz}
\date{Source: arXiv:2411.17317v3 (CC BY 4.0)}
\begin{document}
\maketitle

\noindent\emph{Source and licence.} On combinatorics of plus-one generated line arrangements. Version arXiv:2411.17317v3 (CC BY 4.0), distributed under the Creative Commons Attribution 4.0 International licence, as recorded in the arXiv licence metadata for this exact version and shown on the abstract page https://arxiv.org/abs/2411.17317v3. Full rights notice: licenses/arxiv-2411.17317-CC-BY-4.0.txt.\par\smallskip

\noindent\textit{Editorial scope.} Counted unit: the Thm whose source label is \texttt{comeq} in the article (source0.tex lines 152--157, source label \texttt{comeq}), delivered together with the article's own complete original proof of it (source0.tex lines 159--186, one \texttt{proof} environment of 28 lines). No mathematics is written, repaired or re-derived: statement and proof are the article's own text line for line. Required context retained uncounted: none. Every definition, display and label of those lines is retained unchanged. Omitted from this package and not redistributed: 1--151, 158--158, 187--521. Nothing inside a retained excerpt is omitted or altered: every retained body is delivered exactly as the article's lines are written, so no omission receipt of any kind accompanies this package. Citations: the retained excerpts carry no citation command at all, so no bibliography entry of the article is transcribed and no reference is redistributed. Reference adaptation: none was needed and none was made; every reference inside the delivered unit resolves inside it, so no unresolved reference or citation is produced. Printed slips kept literal: no printed word, symbol, hypothesis or number of the counted statement or of its proof was altered. Layout and class: the article's own class is \texttt{article} with packages including geometry, hyperref, amsfonts, amssymb, amsthm, amsmath, eucal, tabu, url, pgf, array, tikz-cd, pstricks, pstricks-add; the wrapper is the lane's pinned \texttt{amsart} editorial wrapper at 10pt on A4, which restates the theorem-like environments the retained text needs and adds the article's own macro definitions that the retained text needs (none), with extra wrapper packages (bm, bbm, dsfont, xspace, cancel, mathtools). The delivered numbering is the unit's own. Assets: no retained span contains a figure environment, a graphics-inclusion command, a PDF-page inclusion, a table or a TikZ picture; no image file is copied into this package, the package declares no asset dependency and no image file is redistributed with it. Licence: version arXiv:2411.17317v3 of this article is distributed under the Creative Commons Attribution 4.0 International licence, as recorded in the arXiv licence metadata for this exact version and shown on the abstract page https://arxiv.org/abs/2411.17317v3. A full rights notice is delivered at licenses/arxiv-2411.17317-CC-BY-4.0.txt. Characters: every retained excerpt is pure ASCII, so the delivered engine, XeLaTeX with its default Latin Modern text font, reports no missing character.
\begin{thm}
\label{comeq}
     Let $\mathcal{L} \subset \mathbb{P}^{2}_{\mathbb{C}}$ be a plus--one generated arrangement of $d$ lines with the exponents $(d_1, d_2, d_3)$. Then
     $$\sum_{r\geq2}(r-1)\cdot t_r=d_1 d_2+d_3,$$
where $t_r = t_{r}(\mathcal{L})$ is the number of $r$-fold intersection points of $\mathcal{L}$.
\end{thm}

\begin{proof}
    Recall that if $\mathcal{L}$ is line arrangement, then all singularities are quasi--homogeneous and $$\tau(\mathcal{L})=\sum_{r\geq 2}{(r-1)^2\cdot t_r},$$
    where $\tau(\mathcal{L})$ is the total Tjurina number.
    Combining above formula with the naive count, that 
    $$d^2-d=\sum_{r\geq 2}{(r^2-r)\cdot t_r},$$
    we obtain
    $$\tau(\mathcal{L})=\sum_{r\geq 2}(r^2-2r+1)\cdot t_r=\sum_{r\geq 2}(r^2-r)\cdot t_r-\sum_{r\geq 2}(r-1)\cdot t_r=d^2-d-\sum_{r\geq 2}(r-1)\cdot t_r.$$
    Since $\mathcal{L}$ is plus--one generated, we can use Proposition 2.3, and this gives us
    $$\sum_{r\geq 2}(r-1)\cdot t_r=d+d_1d-d_1^2-d_1+d_3-d_2.$$
    The assumption of plus--one generatedness implies $d=d_1+d_2$, and this finally allows us to conclude that 
    $$\sum_{r\geq 2}(r-1)\cdot t_r=d_1+d_2+d_1^2+d_1d_2-d_1^2-d_1+d_3-d_2=d_1d_2+d_3,$$
    which completes the proof.
    

    % Recall that if $\mathcal{L}$ is plus--one generated with exponents $(d_1,d_2,d_3)$, then by Proposition \ref{dimspl} the total Tjurina number of $\mathcal{L}$ can be computed as $$\tau(\mathcal{L}) = (d-1)^{2} - d_{1}(d-d_{1}-1) - (d_{3}-d_{2}+1).$$ Note that all singularities are quasi-homogeneous for line arrangements, so we can use the following count
    % $$\tau(\mathcal{L})=\mu(\mathcal{L})=\sum_{r\geq 2}{(r-1)^2\cdot t_r},$$
    % where $\mu(\mathcal{L})$ is the total Milnor number. Recall also that for line arrangements we have the naive count
    % $$d^2-d=\sum_{r\geq 2}{(r^2-r)\cdot t_r}.$$
    % Thus, we can write $\tau(\mathcal{L})$ as
    % $$\tau(\mathcal{L})=\sum_{r\geq 2}(r^2-2r+1)\cdot t_r=\sum_{r\geq 2}(r^2-r)\cdot t_r-\sum_{r\geq 2}(r-1)\cdot t_r=d^2-d-\sum_{r\geq 2}(r-1)\cdot t_r.$$
    % Therefore, we arrive at the equation
    % $$d^2-2d+1-d_1d+d_1^2+d_1-d_3+d_2-1=d^2-d-\sum_{r\geq 2}(r-1)\cdot t_r,$$
    % and this gives us
    % $$\sum_{r\geq 2}(r-1)\cdot t_r=d+d_1d-d_1^2-d_1+d_3-d_2.$$
    % The assumption that $\mathcal{L}$ is plus--one generated implies $d=d_1+d_2$, and this finally allows us to conclude that 
    % $$\sum_{r\geq 2}(r-1)\cdot t_r=d_1+d_2+d_1^2+d_1d_2-d_1^2-d_1+d_3-d_2=d_1d_2+d_3,$$
    % which completes the proof.
\end{proof}

\end{document}
